\section{Experiments}
\vspace{-1ex}

% \paragraph{Data.} We use the GEOM dataset \citep{GEOM}, consisting of QM9 (average 11 atoms) and DRUGS (average 44 atoms), to train and evaluate our model. 
% We use the same train/val/test molecule splits from \citet{GeoMol} (106586/13323/1000 for
% QM9 and 243473/30433/1000 for DRUGS). 
% The GEOM-DRUGS dataset is also used for the property prediction benchmarks. 
% For the two oracle-based protein docking tasks, we use the Therapeutics Data Commons (TDC)~\citep{tdc}) 9 supported proteins and the 166000 protein-ligand complexes from the CrossDocked~\citep{CrossDocked} training set (see Appendix~\sref{section:protein} for more details), respectively.
% \vspace{-1ex}

We evaluate MCG on RMSD spatial accuracy (\sref{sec:rmsd}), property prediction (\sref{sec:rmsd}), and biochemical quality through flexible (\sref{sec:flex}) and rigid (\sref{sec:rigid}) oracle-based protein docking.
% \paragraph{Baselines.} 
We include the following models for comparison with \modelns, which are previous MCG methods that use the same train/test split: Torsional Diffusion (TD)~\citep{torsional-diffusion}, GeoMol (GM)~\citep{GeoMol}, and when possible, GeoDiff (GD)~\citep{geodiff}.  
% We evaluate all methods over an expanded set of benchmarks to test the robustness of the generated conformers. 
For our model configuration and compute resource breakdown, see Appendix~\sref{section:appendix-breakdown}. 
The details of \modelns's initial RDKit structures, as well as how \model learns to avoid the distribution shift found in TD, can be found in Appendix~\sref{section:mmff}.
% For a discussion of the initial RDKit structures, see Appendix~\sref{section:mmff}.

\vspace{-1ex}
\subsection{GEOM Benchmarks: 3D Coordinate RMSD}
\label{sec:rmsd}
\vspace{-1ex}
We use the GEOM dataset \citep{GEOM}, consisting of QM9 (average 11 atoms) and DRUGS (average 44 atoms), to train and evaluate our model. 
We use the same train/val/test splits from \citet{GeoMol} (QM9: 106586/13323/1000 and DRUGS: 243473/30433/1000).
% The GEOM-DRUGS dataset is also used for the property prediction benchmarks. 
\vspace{-1ex}
\paragraph{Problem setup.}
We report the average minimum RMSD (AMR) between ground truth and generated conformers, and Coverage for Recall and Precision. Coverage is defined as the percentage of conformers with a minimum error under a specified AMR threshold. Recall matches each ground truth conformer to its closest generated structure, and Precision measures the overall spatial accuracy of the generated conformers.
Following \citet{torsional-diffusion}, we generate two times the number of ground truth conformers for each molecule. More formally, for $K=2L$, let $\{C^*_l\}_{l \in [1, L]}$ and $\{C_k\}_{k \in [1, K]}$ respectively be the sets of ground truth and generated conformers:
% \vspace{-4mm}
% \begin{equation}
% \begin{aligned}
% % \text{COV-R} &:= \frac{1}{L}\, \bigg\lvert\{ l \in [1..L]: \exists k \in [1..K], \rmsd(C_k, C^*_l) < \delta \, \bigg\rvert\\
% &\text{COV-Recall} := \frac{1}{L}\, \bigg\lvert\{ l \in [1..L]: \text{min}_{ k\in [1..K]} \rmsd(C_k, C^*_l) < \delta \} \bigg\rvert,\\
% &\text{AMR-Recall} := \frac{1}{L} \sum_{l \in [1..L]}  \text{min}_{ k\in [1..K]} \rmsd(C_k, C^*_l),\\
% % &\text{$f()$} := \{{\min, \mean, \max}\},
% \end{aligned}
% \label{eq:amr}
% \end{equation}
\begin{equation}
\begin{aligned}
% \text{COV-R} &:= \frac{1}{L}\, \bigg\lvert\{ l \in [1..L]: \exists k \in [1..K], \rmsd(C_k, C^*_l) < \delta \, \bigg\rvert\\
&\text{COV-Precision} := \frac{1}{K}\, \bigg\lvert\{ k \in [1..K]: \text{min}_{ l\in [1..L]} \rmsd(C_k, C^*_l) < \delta \} \bigg\rvert,\\
&\text{AMR-Precision} := \frac{1}{K} \sum_{k \in [1..K]}  \text{min}_{ l\in [1..L]} \rmsd(C_k, C^*_l),\\
% &\text{$f()$} := \{{\min, \mean, \max}\},
\end{aligned}
\label{eq:amr}
\end{equation}
% \vspace{-1mm}
where $\delta$ is the coverage threshold. The recall metrics are obtained by swapping ground truth and generated conformers. 
We also report the full RMSD error distributions, as the AMR only gives a small snapshot into overall error behavior.
%as only looking at the minimum score does not yield enough information on the distribution of the generated samples.
% During a true test-time evaluation, the ground truth conformation is unknown. Therefore, we cannot compare it to all $2L$ generated structures to determine the most accurate one for the downstream application.
% Thus, the additions of average (mean) and worst-case (max) evaluation allow us to measure the robustness of the generative model by providing an error bound on each sample. 
% As most downstream tasks that use molecular conformers possess expensive oracle functions, \ie\ protein docking, it is important to be able to generate optimal conformers, as testing all $2L$ samples for each molecule may be infeasible.
% We note that AMR is the central task-agnostic metric. Coverage is derived from AMR and its threshold can be chosen arbitrarily, with its value being task-dependent.

% \paragraph{xTB energy and property prediction.}
% To directly assess the quality of generated samples, we estimate the molecular ensemble properties~\citep{GEOM} over a set of generated conformers. 
% Specifically, similar to TD and GD, we sample min(2K, 32) conformers of 100 GEOM-DRUGS molecules, where K is the number of ground truth conformers. 
% All structures are relaxed with GFN2-xTB~\citep{xtb} with the BGFS optimizer before predicting Energy $E$, HOMO-LUMO Gap $\Delta \epsilon$, and minimum Energy $E_{min}$  in kcal/mol and dipole moment $\mu$ in debye. We report the mean absolute error of generated properties against the ground truth ensemble properties.
% \danny{TODO: not yet have done but Move xtb section to appendix}

\vspace{-1ex}
\paragraph{Advantages and limitations of RMSD-based metrics.}
% The 3D coordinate-based metrics are ideal for evaluating ML methods as they are usually trained with a spatial error-based loss function. However, they are problematic for domain scientists as low spatial error (RMSD) only acts as a proxy for the main quantities they care about, primarily the free energy of the structure~\citep{blau, blau2}. Furthermore, the spatial benchmarks can be separated into two types: precision and recall, both possessing pros and cons. 
3D coordinate-based metrics are commonly used to evaluate ML methods because these models are typically trained using spatial error-based loss functions (\ie,\ MSE). However, for domain scientists, these metrics can be somewhat challenging to interpret. Low spatial error (measured by RMSD) is not directly informative of free energy, the primary quantity of interest to scientists~\citep{blau, blau2}.
Additionally, current spatial benchmarks can be categorized into two distinct types: precision and recall. Each of these metrics comes with its own advantages and limitations.

(1) Precision measures the generation accuracy. It tells us if each generated conformer is close to \textit{any} one of the given ground truth structures, but it does not tell us if we have generated the lowest energy structure, which is the most important at a standardized temperature of 0 K. At industrial temperatures, the full distribution of generated conformers is more important than the ability to generate a single ground truth conformer (for the full RMSD error distribution, see~\fref{fig:rmsd_dist}). 

(2) Recall compares each ground truth to its closest generated conformer. However, in many applications, we are only concerned with obtaining the lowest energy conformer, not all feasible ones. Furthermore, Recall is severely biased by the number of generated conformers for each molecule. Results worsen by up to \textasciitilde{60\%} for all models when we move from the previously set standard sampling budget for each molecule of $2L$ to min(L/32, 1), where $L$ is the number of ground truth conformers (see Appendix ~\fref{fig:2x2grid}). As we sample more molecules, we greatly influence the chance of reducing the AMR-Recall. Because of this dependency on the number of samples, we focus on the Precision metrics, which were consistent across all tested sample size budgets. We note that $L$ is 104 on average for GEOM-DRUGS~\citep{CVGAE}, so even $L$/32 is a reasonable number.
\vspace{-1ex}
\paragraph{Results.}
In~\tref{tab:results_geom}, we outperform all models on QM9, and yield competitive results with TD on DRUGS when using an optimal transport (OT) loss (see Appendix~\sref{section:appendix-qm9}-~\sref{section:appendix-more-drugs} for more details). \model also achieves the lowest overall error distribution, as seen in ~\fref{fig:rmsd_dist}.
\modelns-OT uses an OT loss with the same decoder architecture as in \fref{fig:schematic}, but is no longer autoregressive (see Appendix \eref{eq:ot_loss} for a formal definition).
We also see in Appendix ~\fref{fig:2x2grid} that the recall is heavily dependent on the sampling budget, as performance gaps shrink from \textasciitilde{50\%} to \textasciitilde{5\%}. 
\modelns-OT was trained for 15 hours (2 epochs) compared to TD's 11 days, both on a single A6000. Furthermore, when limited to the same equivariance, \modelns-OT performs predominantly better (Appendix~\tref{tab:results_geom_supp}). 
As \model (autoregressive, without the OT loss) results in a lower overall DRUGS RMSD distribution (\fref{fig:rmsd_dist}), we use this model for the remaining downstream tasks.

\begin{table}[t]
\caption{Quality of ML generated conformer ensembles for the GEOM-QM9 ($\delta=0.5$\AA) and GEOM-DRUGS ($\delta=0.75$\AA) test set in terms of Coverage (\%) and Average RMSD (\AA) Precision. Bolded results are the best, and the underlined results are second best. See Appendix~\sref{section:appendix-qm9}-~\sref{section:appendix-more-drugs} for more details.}
%All models were benchmarked with Torsional Diffusion's evaluation code and retrained if generated molecules were not public (using their public instructions). We note that TD uses higher-order tensor products to maintain equivariance ($\ell = 2$). In contrast, GeoMol, GeoDiff, and \model use scalar-vector operations that are theoretically analogous to $\ell = 1$. \modelns-OT uses an optimal transport (OT) loss with the same decoder architecture as in \fref{fig:schematic}, but is no longer autoregressive. See Appendix \sref{section:appendix-qm9}-\sref{section:appendix-more-drugs} for more details, including the OT equations and recall results. Bolded results are best, and the underlined are second best. \textsuperscript{GeoDiff code would not load, so we were given the GeoDiff generated DRUGS molecules from Torsional Diffusion evaluation on the same test set}\danny{switch to foot}}
\label{tab:results_geom}
\begin{tabular}{lc|cccc|cccc} \toprule
               & \multicolumn{1}{l|}{} & \multicolumn{4}{c|}{QM9-Precision} & \multicolumn{4}{c}{DRUGS-Precision}  \\
                 & \multicolumn{1}{l|}{} & \multicolumn{2}{c}{Coverage $\uparrow$} & \multicolumn{2}{c|}{AMR $\downarrow$} & \multicolumn{2}{c}{Coverage $\uparrow$} & \multicolumn{2}{c}{AMR $\downarrow$} \\
Method  &  & Mean & Med & Mean & Med & Mean & Med & Mean & Med \\ \midrule
\multicolumn{2}{l|}{GeoDiff} &   - & - & -& - & 23.7 & 13.0 & 1.131 & 1.083 \\
\multicolumn{2}{l|}{GeoMol} &   75.9 & \textbf{100.0} & 0.262  & 0.233 & 40.5 & 33.5 & 0.919 & 0.842 \\
\multicolumn{2}{l|}{Torsional Diffusion ($\ell = 2$)} & \underline{78.4} & \textbf{100.0} &  \underline{0.222}  &  \underline{0.197} & \textbf{52.1} & \textbf{53.7} & \textbf{0.770} & 0.720 \\
\multicolumn{2}{l|}{\modelns-OT} & \textbf{80.2} & \textbf{100.0} & \textbf{0.149} & \textbf{0.107} & \underline{52.0} & \underline{52.1} & \underline{0.836} & \textbf{0.694} \\
\bottomrule
\end{tabular}
\vspace{-3ex}
\end{table}

\begin{figure}[hbt]%[H] % You can use [h] to try to place the figure here
  \vspace{-3ex}
  \begin{minipage}{0.5\textwidth}
    \centering
    \includegraphics[width=\textwidth]{figures/rmsd2.png}
    % \caption{GEOM-DRUGS test set RMSD distributions for each tested ML method. Here we provide the raw distribution that is used to create the precision and recall metrics.}
    % \caption{GEOM-DRUGS test set RMSD \\ distributions for each ML model.}
    \vspace{-10pt}
    % \vspace{-5pt}
  \end{minipage}
  \begin{minipage}{0.5\textwidth}
    \centering
    \captionof{table}{Property prediction: Mean absolute error of generated vs. ground truth ensemble properties for $E$, HOMO-LUMO gap $\Delta \epsilon$, $E_{min}$ (kcal/mol), and dipole moment $\mu$ (debye).
    }
%For 100 molecules, min(2L, 32) where L is the number of ground truth conformers are generated and relaxed with GFN2-xTB with the BGFS optimizer before predicting Energy $E$, HOMO-LUMO Gap $\Delta \epsilon$, and minimum Energy $E_{min}$  in kcal/mol and dipole moment $\mu$ in debye via xTB. We report the mean absolute error of generated v.s. ground truth ensemble properties.}
    \begin{tabular}{l|llll} 
        \toprule
        & $E$ & $\mu$ & $\Delta \epsilon$ & $E_{min}$ \\
        \midrule
        GeoMol & 28.80 & 1.475 & 4.186 & 0.267 \\
        Torsional Diffusion & 16.75 & 1.333 & 2.908 & 0.096 \\
        \model & \textbf{12.41} & \textbf{1.250} & \textbf{2.522} & \textbf{0.049} \\
        \bottomrule
        \end{tabular}
    % \caption{Table caption}
    \label{tab:xtb_energy}
  \end{minipage}
   \caption{GEOM-DRUGS test set RMSD error distributions for each ML model. %Here, we provide the raw distribution that is used to calculate the precision and recall metrics.
     }
 \label{fig:rmsd_dist}
 \vspace{-3ex}
\end{figure}
% \vspace{-1ex}
\subsection{GEOM Benchmarks: Property Prediction}
\label{sec:prop}
% We also use the GEOM-DRUGS dataset to evaluate property prediction benchmarks.
\vspace{-2ex}
\paragraph{Problem setup.} %For 100 molecules from GEOM-DRUGS, min(2L, 32) where L is the number of ground truth conformers are generated and relaxed with GFN2-xTB with the BGFS optimizer before predicting Energy $E$, HOMO-LUMO Gap $\Delta \epsilon$, and minimum Energy $E_{min}$  in kcal/mol and dipole moment $\mu$ in debye via xTB. We report the mean absolute error of generated v.s. ground truth ensemble properties.
% For further discussions about RMSD and energy-based benchmarks, see Appendix~\sref{section:geom-xtb}).
We generate and relax min(2L, 32) conformers (L ground truth) for 100 molecules from GEOM-DRUGS using GFN2-xTB with the BFGS optimizer. We then predict various properties, including Energy ($E$), HOMO-LUMO Gap ($\Delta \epsilon$), minimum Energy ($E_{min}$) in kcal/mol, and dipole moment ($\mu$) in debye via xTB~\citep{xtb}. The mean absolute error of the generated ensemble properties compared to ground truth is reported. 
% Further discussions can be found in Appendix~\sref{section:geom-xtb}.

\vspace{-2ex}
\paragraph{Results.}\tref{tab:xtb_energy} demonstrates \modelns's ability to generate the lowest energy structures with the most accurate chemical properties. For further discussions, please see Appendix~\sref{section:geom-xtb}.


\vspace{-1ex}
\subsection{Flexible Oracle-based Protein Docking}
\label{sec:flex}
We evaluate MCG models, pretrained on GEOM-DRUGS, using nine protein docking oracle functions provided by the Therapeutics Data Commons (TDC)~\citep{tdc}.
% We evaluate \model pretrained on GEOM-DRUGS on flexible oracle-based protein docking. We use the Therapeutics Data Commons (TDC)~\citep{tdc}) 9 supported protein docking oracle functions.
% \ak{I think we should just make one section for benchmarks, because then we can also add the property prediction one}
% To promote more comprehensive evaluations of generated conformers, we use oracle-based protein docking as a meaningful downstream task to evaluate the binding affinity of the generated conformers, which can be thought of as dynamic property prediction \ak{We may want to be careful to use ``dynamic property prediction'' when it isn't really defined}. 
% Since current benchmarks mainly focus on gauging the effectiveness of the machine learning objective and less on the chemical feasibility and downstream use of the generated conformers, 
\vspace{-1ex}
\paragraph{Problem setup.} Starting with a known 2D ligand\footnote{No ground truth 3D structures. All ligand SMILES taken from Protein Data Bank: https://www.rcsb.org/} molecule, protein, and desired 3D protein binding pocket, we measure conformer quality by comparing the predicted binding affinity of generated conformers of each MCG method.
% Here we measure conformer quality by comparing the predicted binding affinity of generated conformers for known 2D ligand molecules for known proteins and docking locations.
TDC's protein docking oracle functions take in a ligand SMILES string, generate a 3D conformer, and try multiple poses, before ultimately returning the best binding affinity via Autodock Vina's flexible docking simulation.
We augment TDC with the ability to query ML models pretrained on GEOM-DRUGS, instead of the built-in RDKit + MMFF approach for MCG.
% We use the standardized protein docking oracle functions of TDC and augment them with the ability to query ML models pretrained on GEOM-DRUGS, instead of the built-in RDKit + MMFF approach. We use this to generate and evaluate conformers of known ligands.
% TDC uses Autodock Vina's flexible docking simulation to optimize the input ligand conformer and predict the binding affinity.
For each evaluated MCG method, we generate 50 conformers for each of the nine ligands and report the best (lowest) binding affinity.
Given that the ligand and protein identity and the protein pocket are fixed, this task measures the quality of 3D conformer coordinates through their binding efficacy to the specified pocket. We note that this task is indicative of real-world simulation workflows.


\begin{table}[b]%[H]
\vspace{-1ex}
\begin{center}
% \centering
\caption{Quality of best generated conformer for known protein ligands for all 9 proteins from the TDC library. Quality is measured by free energy change (kcal/mol) of the binding process with AutoDock Vina's flexible docking simulation ($\downarrow$ is better).
}
%where low energy (negative values) is more favorable. }
\label{tab:results_oracle}
% \hspace{-10mm}
\begin{tabular}{l|ccccccccc}
\toprule
                    & \multicolumn{9}{c}{Best Protein-Conformer Binding Affinity ($\downarrow$ is better)}                                                                                                                   \\
Method              & 3PBL           & 2RGP            & 1IEP            & 3EML           & 3NY8            & 4RLU & 4UNN            & 5M04            & \multicolumn{1}{l}{7L11}          \\ \hline
RDKit + MMFF       & -8.26          & -11.42          & -10.75          & -9.26          & -9.69           &   -8.72   & -9.73           & -9.53           & -9.19                             \\
GeoMol              & -8.23          & -11.49          & -11.16          & -9.39          & \textbf{-11.66} &  -8.85    & -10.28          & -9.31           & -9.29                             \\
Torsional Diffusion & -8.53          & -11.34          & -10.76          & -9.25          & -10.32          &   -8.96   & -10.65          & -9.61           & -9.10                             \\
CoarsenConf       & \textbf{-8.81} & \textbf{-12.93} & \textbf{-16.43} & \textbf{-9.82} & -11.26          &   \textbf{-9.54}   & \textbf{-11.62} & \textbf{-14.00} & \textbf{-9.43}                    \\ \hline
% CC-next best ML model $\Delta$         & \textbf{-0.28}  & \textbf{-1.44}   & \textbf{-5.27}   & \textbf{-0.43}  & 0.40           &   \textbf{-0.58}   & \textbf{-0.97}   & \textbf{-4.39}   & \multicolumn{1}{l}{\textbf{-0.14}} \\
% CC-TD $\Delta$             & \textbf{-0.28}  & \textbf{-1.59}   & \textbf{-5.67}   & \textbf{-0.57}  & \textbf{-0.94}   &  \textbf{-0.58}    & \textbf{-0.97}   & \textbf{-4.39}   & \multicolumn{1}{l}{\textbf{-0.33}}
\end{tabular}
\end{center}
\vspace{-2ex}
\end{table}
\vspace{-1ex}
\paragraph{Results.}
\model significantly outperforms prior MCG methods on the TDC oracle-based affinity prediction task (~\tref{tab:results_oracle}).
\model generates the best ligand conformers for 8/9 tested proteins, with improvements of up to 53\% compared to the next best method. 
\model is 1.46 kcal/mol better than all methods when averaged over all 9 proteins, which corresponds to a 14.4\% improvement on average compared to the next best method.
% \vspace{-1ex}
\subsection{Rigid Oracle-based Protein Docking}
\vspace{-1ex}
\label{sec:rigid}
We evaluate MCG models on rigid oracle-based protein docking. We use the 166000 protein-ligand complexes from the CrossDocked~\citep{CrossDocked} training set (Appendix~\sref{section:protein} for more details).

\vspace{-1ex}
\paragraph{Problem setup.} Similar to the flexible docking task (\sref{sec:flex}), we can generate conformers of known ligands for known protein pockets, but now have them only undergo a rigid pocket-specific energy minimization before predicting the binding affinity. 
% Our evaluation uses the CrossDocked trainset, which has tens of millions of conformer-protein complexes at our disposal (See dataset details in Appendix~\sref{section:protein}).
% The CrossDocked dataset is derived from PDBBind but uses smina~\citep{smina}, a derivative of AutoDock Vina with more explicit scoring control, to generate the protein-conditioned ligand structures to yield ground truth data. We note that based on the raw data, 2.2 billion conformer-protein interactions are possible, but we filtered out any ground-truth example that AutoDock Vina failed to score. 
To handle the fact that MCG generates ligand structures in a vacuum and has no information about the target protein, we assume the centroid of the ground truth ligand is accurate, and translate each generated structure to match the center.
% Furthermore, in the oracle-based task, each ligand is known to fit well in the given protein. CrossDocked, on the other hand, is comprised of various ligand-protein interactions, not all of which are optimal.

We use AutoDock Vina~\citep{vina} and its local BFGS energy optimization (similar to that done with xTB), \ie\ a relaxation of all tested structures (including the ground truth), following the SBDD benchmarking framework~\cite{targetdiff}. We report the difference between the generated structure’s minimum binding affinity and that of the ground truth 3D ligand. 
% The Vina minimization refines the input ligand's position and orientation to find the most energetically favorable configuration without changing the internal torsion angles of the ligand. 
% Compared to the more rigorous search of the flexible docking simulation, minimization is used for structural refinement and energy minimization of a conformer in isolation. 
Unlike the docking simulation, Vina's energy minimization does not directly adjust the torsion angles or internal degrees of freedom within the ligand. Instead, it explores different atomic positions and orientations of the entire ligand molecule within the binding site of the protein to find energetically favorable binding poses. 
By further isolating the MCG-generated structures, this task better evaluates the generative capacity of MCG models. % compared to the prior docking task.
% We emphasize that all tested MCG-generated structures undergo the same structure optimization, even the ground truth data, to create an equal and fair conformer benchmark.
% Following ~\cite{targetdiff} better generated outputs (in our case conformers) will result in a more negative binding affinity, aiming to match or exceed the ground truth data as much as possible. 
While the original SBDD task~\citep{pocket2mol, targetdiff} reports the Vina minimization score as its main metric, it requires the 3D ligand to be generated from scratch. Here, we use the SBDD framework to isolate the generated ligand conformer as the only source of variability to evaluate the biological and structural quality of MCG models in an unbiased fashion. 
% We reiterate that by isolating the only source of variability to the MCG-generated conformer, we can use the desirable task of ligand-protein binding to score the generative ability of MCG models in an unbiased fashion 
% \ak{This is redundant from above, I think we can merge sentences}. 

% We note that while it takes on the orders of hours to generate 1.2M conformers (100 conformers per molecule), it takes on the orders of \textasciitilde{weeks to months} to score each conformer for up to the 2,358 unique proteins for each evaluated method (evaluation time is 100x the time to score the ground truth data as we generate 100 conformers per molecule). As a result, we report the results for the first 100,000 conformer-protein interactions.
\vspace{-1ex}
\paragraph{Results.} We report the results for 100,000 unique conformer-protein interactions; note that there is a large cost to run the binding affinity prediction (see Appendix~\sref{section:protein} for more details).
We also emphasize that the presented evaluation is not to be confused with actual docking solutions, as a low-energy conformer is not always guaranteed to be the best binding pose. Instead, we employ an unbiased procedure to present empirical evidence for how \model can generate input structures to Vina that significantly outperform prior MCG models in achieving the best binding affinities.

% In this regard, we use Vina's predicted binding affinity as a dynamic property prediction task.
% \vspace{-1ex}
% \paragraph{Data-constrained setting.} \danny{replace this with explanation of the 4 benchmarks pointing to supplementary for further discussion number each} In real-world applications like polymer design, the availability of data is frequently limited, and accompanied by a scarcity of conformers for each molecule. The current datasets, QM9 and DRUGS, do not mimic this setting very well. For example, on average QM9 has 15 conformers per molecule and DRUGS has 104 per molecule---both datasets have significantly more conformers than in an experimental drug design setting. Given this, rather than training on the first 30 conformers as done in TD, we train on the first five and two conformers (typically those with the largest Boltzmann weight) for QM9 and DRUGS, respectively. This corresponds to 41\% of TD's training set or 23\% of the overall dataset for QM9, and 7.3\% of TD's training set or 1.7\% of DRUGS overall.
% Even with less data, we outperform models trained on significantly larger training sets.
% \vspace{-1ex}
% \paragraph{Results.}
\fref{fig:results_sbdd} further demonstrates \modelns's superior performance on orders of magnitude more protein complexes than the prior flexible oracle task. 
\model decreases the average error by 56\% compared to TD, and is the only method not to exhibit bimodal behavior with error greater than zero. 
% We hypothesize as none of the evaluated MCG methods generate conformers conditioned on the given protein pockets, the ability to meet and even surpass ground truth structures stems from the protein pocket complexity.
% More simply, open pockets make it easier for AutoDock Vina to optimize the given conformer, but the starting MCG-generated position remains extremely important.
% Furthermore, if the initial generated structure did not matter, the distributions for each MCG model would be identical.
We hypothesize that the success of MCG methods in matching ground truth structures is influenced by the complexity of protein pockets. In simpler terms, open pockets better facilitate Vina's optimization, but the initial position generated by MCG remains crucial.
Furthermore, if the initial MCG-generated structure did not matter, the distributions for each MCG model would be identical.
We also note that the lowest energy conformer is not always the optimal ligand structure for binding, but in our experiments, it yields the best input for the Vina-based oracles. 
% In this instance, all ligand conformers (generated and ground truth) undergo the same energy minimization prior to affinity scoring.
Overall, \model best approximates the ground truth ligand conformers of CrossDocked and generates the best structures for Vina's rigid energy relaxation and binding affinity prediction.
\vspace{-1ex}
\begin{figure}[t]%[H] % You can use [h] to try to place the figure here
  \vspace{-3ex}
  \begin{minipage}{0.5\textwidth}
    \centering
    \includegraphics[width=\textwidth]{figures/sbdd.png}
    % \vspace{-10pt}
    \label{fig:sbdd}
    % \vspace{-5pt}
    \vspace{-3ex}
  \end{minipage}
  \begin{minipage}{0.5\textwidth}
    \centering
    \captionof{table}{Binding affinity error distribution statistics in kcal/mol (more negative is better).}
    \begin{tabular}{c|cccc}
      \toprule
      Method & Mean & Min \\
      \hline
      GeoMol & 2.476 & -8.523 \\
      Torsional Diffusion & 1.178 & -6.876  \\
      \model & \textbf{0.368} & \textbf{-8.602} \\
      \hline
    \end{tabular}
    % \caption{Table caption}
    \vspace{-3ex}
  \end{minipage}
   \caption{ Binding affinity ($\downarrow$ is better) error distributions for 100k conformer-protein complexes in the CrossDocked dataset. The error is the difference in binding affinity between the generated and ground truth energy minimized 3D ligands.
   }
   %We present binding affinity error distributions for the first 100k conformer-protein complexes in the CrossDocked dataset, following ~\cite{targetdiff}. The error is calculated as the difference in binding affinity between the generated (via MCG) and ground truth energy minimized 3D ligands. AutoDock Vina's energy minimization optimizes the initial 3D structure (via roto translations) to fit the binding pocket while preserving internal geometry.}
     \label{fig:results_sbdd}
     \vspace{-3ex}
\end{figure}


% \vspace{-2ex}


% \vspace{-2ex}
% \begin{figure}[h]%[H] % You can use [h] to try to place the figure here
%   \vspace{-3ex}
%   \begin{minipage}{0.5\textwidth}
%     \centering
%     \includegraphics[width=\textwidth]{figures/sbdd.png}
%     % \vspace{-10pt}
%     \label{fig:sbdd}
%     % \vspace{-5pt}
%     \vspace{-3ex}
%   \end{minipage}
%   \begin{minipage}{0.5\textwidth}
%     \centering
%     \captionof{table}{Binding affinity error distribution statistics in kcal/mol (more negative is better).}
%     \begin{tabular}{c|cccc}
%       \toprule
%       Method & Mean & Min \\
%       \hline
%       GeoMol & 2.476 & -8.523 \\
%       Torsional Diffusion & 1.178 & -6.876  \\
%       \model & \textbf{0.368} & \textbf{-8.602} \\
%       \hline
%     \end{tabular}
%     % \caption{Table caption}
%     \vspace{-3ex}
%   \end{minipage}
%    \caption{Stucture-based Drug Discovery (SBDD) protein-ligand binding affinity evaluation. Starting from CrossDocked, we take the 11,735 unique ligands, generate 100 conformers, and score each for all of the potential 2,358 proteins. The scoring process uses AutoDock Vina and its local BFGS energy optimization, \ie a relaxation of all tested structures, including ground truth following ~\cite{targetdiff}, and reporting the difference between the generated structure's minimum binding affinity and that of the minimized ground truth 3D ligand. We report the error distribution and accompanying statistics of the first 100,000 conformer-protein complexes.
%    % \danny{show example of poor and good explain how to get the center point so no translation and how autodock cvina how it handles the fitting. more intricate pockets result in needing more optimization vina clear what we did vs what we took from other code. instead of conformer in vacuum we can apply to pockets or protein information it has never seen see how crossdocked was generated. so this is unconditional generation applied towards a conditional task to yield an unbiased evaluation}
%    }
%      \label{fig:results_sbdd}
%      \vspace{-3ex}
% \end{figure}


% \begin{figure}[t]
% \centering
%     \includegraphics[width=0.5\textwidth]{figures/sbdd.png}
%     \vspace{-10pt}
%     \caption{Stucture-based Drug Discovery (SBDD) protein-ligand binding affinity evaluation. For 166,500 ground truth binding interactions from the train set of the CrossDocked dataset comprised of 11,735 unique ligands and 2,358 unique proteins, we generate 100 conformers per ligand resulting in almost 1.2 million structures for each MCG method. We use AutoDock Vina and its minimization or relaxation on all tested structures, including ground truth following \danny{cite target diff}, and report the difference between the generated structure's minimum binding affinity and that of the ground truth 3D ligand. Better 3D conformers will result in a more negative binding affinity, aiming to match or exceed the ground truth data as much as possible.
%      }
%     \label{fig:sbdd}
%     \vspace{-5pt}
% \end{figure}

% \subsection{Metrics}
% \vspace{-1ex}
% \paragraph{Metrics.} We report both RMSD (AMR) and Coverage for Recall (R) and Precision (P). Recall measures the number of ground truth conformers that are correctly predicted, and Precision measures the accuracy of the generated conformers.
% Following \citet{torsional-diffusion}, we generate two times the number of ground truth conformers for each molecule. More formally, for $K=2L$ let $\{C^*_l\}_{l \in [1, L]}$ and $\{C_k\}_{k \in [1, K]}$ be respectively the sets of ground truth and generated conformers:
% % \vspace{-4mm}
% \begin{equation}
% \begin{aligned}
% % \text{COV-R} &:= \frac{1}{L}\, \bigg\lvert\{ l \in [1..L]: \exists k \in [1..K], \rmsd(C_k, C^*_l) < \delta \, \bigg\rvert\\
% &\text{COV($f$)-R} := \frac{1}{L}\, \bigg\lvert\{ l \in [1..L]: f_{ k\in [1..K]} \rmsd(C_k, C^*_l) < \delta \} \bigg\rvert,\\
% &\text{AMR($f$)-R} := \frac{1}{L} \sum_{l \in [1..L]}  f_{ k\in [1..K]} \rmsd(C_k, C^*_l),\\
% &\text{$f()$} := \{{\min, \mean, \max}\},
% \end{aligned}
% \end{equation}
% % \vspace{-1mm}
% where $\delta$ is the coverage threshold. The precision metrics are obtained by swapping ground truth and generated conformers. 
% During a true test-time evaluation, the ground truth conformation is unknown. Therefore, we cannot compare it to all $2L$ generated structures to determine the most accurate one for the downstream application.
% Thus, the additions of average (mean) and worst-case (max) evaluation allow us to measure the robustness of the generative model by providing an error bound on each sample. 
% As most downstream tasks that use molecular conformers possess expensive oracle functions, \ie\ protein docking, it is important to be able to generate optimal conformers, as testing all $2L$ samples for each molecule may be infeasible.
% We note that AMR is the central task-agnostic metric. Coverage is derived from AMR and its threshold can be chosen arbitrarily, with its value being task-dependent.

\vspace{-1ex}
% \subsection{Results}
% \label{section:results}
% \vspace{-1ex}

% \paragraph{GEOM}
% Our results, shown in~\tref{tab:results_geom} and ~\fref{fig:rmsd_dist} demonstrate the best precision and lower overall error distribution. \model outperforms GM and GD while maintaining competitive results and even exceeding TD for QM9 and DRUGS. We also see in Appendix ~\fref{fig:2x2grid} how the recall is heavily dependent on the sampling budget as performance gaps shrink from \textasciitilde{50\%} to \textasciitilde{5\%}. Given the small margin between each method, we turn to energy-based tasks that measure the chemical similarity between the generated and ground truth structures. \tref{tab:xtb_energy} demonstrates \modelns's ability to generate the lowest energy structures with the most accurate chemical properties. We note \modelns-OT $\dagger$ was trained for 15 hours (2 epochs) compared to TD's 11 days both on a single A6000\ak{14 days on what GPU?}. Furthermore, when limited to the same equivariance, \model performs best (Appendix~\sref{section:appendix-more-drugs}). As \model with an autoregressive decoder trained without the OT loss results in a lower overall DRUGS RMSD distribution, we use that model for the remaining downstream task evaluations.
% \vspace{-1ex}
% \paragraph{Downstream binding affinity}
% \model significantly outperforms prior MCG methods on the TDC oracle-based affinity prediction task (~\tref{tab:results_oracle}).
% \model generates the best ligand conformers for 8/9 tested proteins, with improvements of up to 53\% compared to the next best method. 
% \model is 1.46 kcal/mol better than all methods when averaged over all 9 proteins, which corresponds to a 14.4\% improvement on average compared to the next best method.

% \fref{fig:results_sbdd} further demonstrates \modelns's superior performance on orders of magnitude more protein complexes than the prior oracle-based task. 
% \model decreases the average error by 56\% compared to TD and is also able to outperform the ground truth data in almost half of its tested samples. 
% We hypothesize as none of the evaluated MCG methods generate conformers conditioned on the given protein pockets, the ability to meet and even surpass ground truth structures stems from the protein pocket complexity. More simply, open pockets make it easier for AutoDock Vina to optimize the given conformer, but the starting MCG-generated position remains extremely important. 
% We also note that the lowest energy conformer is not always the optimal ligand structure for binding, but in our experiments, we see it does yield the best input for the Vina-based oracles. 
% % In this instance, all ligand conformers (generated and ground truth) undergo the same energy minimization prior to affinity scoring.
% Overall \model best approximates the ground truth ligand conformers of CrossDocked and generates the best structures for Vina's rigid energy optimization and binding affinity prediction.
% % \paragraph{QM9.} Our results, shown in~\tref{tab:results_qm9}, demonstrate a significant improvement over both RDKit+MMFF and Torsional Diffusion (TD) on the mean and max metrics.
% % We have effectively shown that our model can improve upon RDKit initializations, resulting in state-of-the-art generalization. 
% % \model does so using 59\% of TD's data and for only 5 epochs compared to 250.

% % \vspace{-1ex}
% % \paragraph{DRUGS.} \tref{tab:results_drugs} demonstrates \modelns's ability to best generate robust low-error conformations for large drug-like molecules. With even further constraints on the data (only two conformers per molecule) and total iterations (5 epochs), \model achieves superior AMR results for the more comprehensive mean and max metrics. For comparison, TD uses 30 conformers for 250 epochs.
% % We also demonstrate that \model outperforms GeoDiff~\citep{geodiff}, a diffusion model over 3D coordinate point clouds, and GeoMol, which learns to predict a minimal set of geometric quantities ($\ie$  angles
% % and bond lengths) and then reconstruct the local and global structures of the conformation.

% \vspace{-1ex}
% \paragraph{Takeaway \danny{may bump}.}
% Both \model and TD optimize $p(X|\gR)$ but utilize the RDKit approximations $\gR$ in different ways. TD learns to update the torsion angles of $\gR$ while \model leverages CG information to inform geometric updates (coordinates, distances, and torsion angles) to translate $\gR$ to $X$.
% Unlike TD, which uses a preprocessing optimization procedure to generate substitute ground truth conformers that mimic $p(\gR)$, 
% \model directly learns from both $X$ and $\gR$ through its hierarchical graph matching procedure. Thus directly addressing the distributional shift problem.
% We hypothesize that this, along with our angle-based CG strategy, leads to our observed improvements.
% Overall, \model provides a comprehensive framework for accurate conformer generation that can be directly applied for downstream tasks such as oracle-based protein docking. 
% Additional results are in Appendix \sref{section:appendix-more-drugs}.
% \vspace{-2ex}

